problem:  
Let \((M^n, g(t))_{t\in(-\infty,0]}\) be a complete ancient solution to the Ricci flow with bounded nonnegative curvature operator and assume it is κ-noncollapsed on all scales for some \(\kappa>0\).  
For such a flow, define the **asymptotic tangent flow at infinity** as any pointed limit of the rescaled metrics  
\[
g_k(s) := |t_k|^{-1} g(t_k + |t_k|s),\qquad s\in(-\infty,0],
\]
for a sequence \(t_k\to-\infty\). Suppose that for one such sequence, the limit \((M_\infty,g_\infty(s),x_\infty)\) is a nonflat gradient shrinking soliton.  

Now consider the collection of all possible blow-down limits obtained from *different* sequences \(t_k\to -\infty\). It is not known in general whether uniqueness of blow-down limits (up to isometry and scaling) must hold.  

**Open problem:**  
Does there exist a κ-noncollapsed ancient Ricci flow with nonnegative curvature operator for which **two non-isometric gradient shrinking solitons** arise as distinct blow-down limits at \(t\to -\infty\)?  
In other words, can different asymptotic tangent flows coexist in the distant past, or must all backward limits of such an ancient flow be mutually isometric shrinkers?  

Equivalently, is the asymptotic shrinker at \(t=-\infty\) *unique* (up to scaling and isometry) among nonnegative-curvature ancient κ-solutions?

---

Why is it a "good" problem:  
This problem targets one of the most subtle and foundational uncertainties in the classification of ancient Ricci flows—**the uniqueness and stability of asymptotic tangent flows**. It is mathematically well-defined, yet remains unresolved even in dimension four.  

1. **Conceptual simplicity:**  
   The question is elementary to state: “Can an ancient κ-solution have more than one distinct backward soliton limit?” Yet the dynamics of such backward asymptotics probe the very structure of high-dimensional Ricci flow evolution.  

2. **Bridge between analysis and geometry:**  
   The problem sits at the interface of **parabolic analysis**, **geometric limit theory**, and **dynamical systems interpretation** of the Ricci flow. Blow-down uniqueness would assert that ancient κ-solutions are asymptotically governed by a single soliton “attractor,” mirroring the uniqueness of tangent cones in geometric measure theory.  

3. **Analytic and geometric difficulty:**  
   Proving or disproving uniqueness would require new control over global monotone quantities (such as Perelman’s entropy) and their limit behavior, or constructing analytic counterexamples mimicking multiple attractors. Both directions are profoundly deep and potentially paradigm-shifting.  

4. **Potential implications:**  
   A positive answer—uniqueness—would suggest a universal “backward dynamical rigidity” of Ricci flow and might yield new classification results for ancient solutions. A counterexample would uncover entirely new geometric phenomena: ancient flows with multiple distinct asymptotic geometric phases, expanding the landscape of possible Ricci evolutions.  

Thus, this problem possesses simplicity in statement, depth in meaning, and connects analytic rigidity, geometric limit theory, and dynamical stability—making it an exemplary "good" research-level problem.